% Vietnamese translation for OI Public Library
% Source-PDF: IOI中国国家候选队论文/2002/骆骥.pdf
\documentclass[11pt]{article}
\usepackage{oipl-vn}

\newcommand{\Oh}{\mathcal{O}}
\newcommand{\Win}{\mathsf{W}}
\newcommand{\Lose}{\mathsf{L}}

\title{Hai hướng tư duy khi giải bài toán đối sách\\
\large Từ bài toán lấy sỏi}
\author{Luo Ji}
\date{}

\begin{document}
\maketitle

\origtitle{Two approaches to solving game-strategy problems: starting from the stone-taking problem}
\sourcepath{IOI China candidate team papers / 2002 / Luo Ji.pdf}

\begin{abstract}
Vận trù học là một ngành còn trẻ, bao gồm quy hoạch, lý thuyết đồ thị,
lý thuyết đối sách, lý thuyết hàng đợi, lý thuyết lưu trữ, v.v. Trong các
kỳ thi tin học, dạng đối sách thường gặp nhất là trò chơi hai người, trong đó
khi đối thủ luôn chọn chiến lược tối ưu thì bên ta hoặc có chiến lược thắng
chắc, hoặc buộc thua. Bài viết bắt đầu từ một bài lấy sỏi cổ điển, qua đó
trình bày hai hướng tư duy cơ bản: phương pháp tổng quát dựa trên đồ thị
chuyển trạng thái, và phương pháp đặc thù dựa trên việc tìm trạng thái cân
bằng cùng quy luật ra quyết định.
\end{abstract}

\tableofcontents

\section{Bài toán lấy sỏi}

Có \(N\) viên sỏi. Hai người chơi Giáp và Ất luân phiên lấy sỏi. Mỗi lần phải
lấy ít nhất một viên, và được lấy nhiều nhất bằng hai lần số sỏi mà đối thủ
vừa lấy ở lượt trước. Giáp đi trước; ở lượt đầu, Giáp có thể lấy tùy ý số
lượng sỏi, nhưng không được lấy hết. Người đầu tiên không còn sỏi để lấy là
người thua.

Câu hỏi là: Giáp có thể thắng hay không? Trong bản bài toán trên slide,
\(1<N<100\).

\subsection{Trạng thái}

Trong bài toán này, kết quả phụ thuộc vào hai yếu tố chính:
\begin{itemize}
  \item số sỏi hiện còn, ký hiệu là \(N\);
  \item số sỏi tối đa được lấy trong lượt hiện tại, ký hiệu là \(K\).
\end{itemize}

Ta dùng cặp \((N,K)\) để mô tả một trạng thái. Trạng thái ban đầu là
\((N,N-1)\), vì Giáp được lấy từ \(1\) đến \(N-1\) viên. Nếu ở trạng thái
\((N,K)\) người chơi lấy \(x\) viên, với
\[
  1\le x\le \min(N,K),
\]
thì trạng thái dành cho người kế tiếp là
\[
  (N-x,2x).
\]
Khi \(N=0\), người đến lượt không có nước đi, nên trạng thái đó là trạng thái
thua cho người sắp đi.

\subsection{Ví dụ N bằng 4}

Với \(N=4\), trạng thái ban đầu là \((4,3)\). Cây chuyển trạng thái đi từ
trên xuống như sau:
\[
\begin{array}{c}
(4,3)\\[0.4em]
\begin{array}{ccc}
\xrightarrow{x=1} & \xrightarrow{x=2} & \xrightarrow{x=3}\\[-0.2em]
(3,2) & (2,4) & (1,6)
\end{array}
\end{array}
\]
Vì nếu \(K\ge N\) thì người đang đi có thể lấy hết sỏi, các trạng thái
\((2,4)\) và \((1,6)\) là trạng thái thắng. Nhánh còn lại có dạng
\[
\begin{array}{rcl}
(3,2) &\to& (2,2)\quad\text{hoặc}\quad (1,4),\\
(2,2) &\to& (1,2)\quad\text{hoặc}\quad (0,4),\\
(1,2) &\to& (0,2).
\end{array}
\]
Từ đáy cây suy ngược lên, ta có:
\[
\begin{array}{c|c|l}
\text{Trạng thái} & \text{Loại} & \text{Lý do}\\ \hline
(0,0) & \Lose & \text{không có nước đi}\\
(1,K),\ K\ge 1 & \Win & \text{lấy hết một viên}\\
(2,2) & \Win & \text{có thể đi tới }(0,4)\\
(3,2) & \Lose & \text{mọi nước đi đều tới trạng thái thắng}\\
(4,3) & \Win & \text{lấy một viên, để lại }(3,2)
\end{array}
\]
Ở đây \(\Win\) nghĩa là người sắp đi thắng chắc, còn \(\Lose\) nghĩa là người
sắp đi thua chắc.

\section{Hướng thứ nhất: phương pháp tổng quát}

Phương pháp tổng quát bắt đầu từ trạng thái ban đầu, xét từ trên xuống tất cả
trạng thái có thể xuất hiện, rồi dựng cấu trúc topo của các phép chuyển trạng
thái. Điểm mạnh của hướng này là có quy tắc thắng thua và cách cài đặt khá
thống nhất, nên phạm vi áp dụng rộng. Những bài như bài lấy số của IOI 1996
hoặc \textit{Ioiwari} của IOI 2001 đều có thể xử lý theo hướng này.

\subsection{Quy tắc thắng thua}

Thắng thua của một trạng thái được quyết định bởi thắng thua của các trạng
thái con trực tiếp.

\begin{itemize}
  \item Nếu một trạng thái không có trạng thái con, tức là đã kết thúc ván,
  thì thắng thua được quyết định trực tiếp theo điều kiện đề bài.
  \item Nếu một trạng thái có ít nhất một trạng thái con là trạng thái thua
  cho người sắp đi, thì trạng thái hiện tại là trạng thái thắng. Người chơi
  chỉ cần đi tới trạng thái con đó.
  \item Nếu mọi trạng thái con đều là trạng thái thắng cho người sắp đi, thì
  trạng thái hiện tại là trạng thái thua.
\end{itemize}

Với bài lấy sỏi, đặt \(f(N,K)\in\{\Win,\Lose\}\). Ta có công thức quy hoạch:
\[
f(N,K)=
\begin{cases}
\Lose, & N=0,\\
\Win, & \exists x,\ 1\le x\le \min(N,K),\
        f(N-x,2x)=\Lose,\\
\Lose, & \text{ngược lại.}
\end{cases}
\]
Đáp án của bài toán ban đầu là \(f(N,N-1)\).

\subsection{Quy tắc mở rộng}

Trong một số bài, ngoài thắng thua, ta còn cần ghi lại lợi ích lớn nhất hoặc
nhỏ nhất mà người đi trước có thể bảo đảm khi thắng hoặc khi thua. Khi đó
trạng thái không chỉ mang nhãn \(\Win/\Lose\), mà còn mang một giá trị số.
Quy tắc chuyển tiếp được điều chỉnh theo mục tiêu của đề bài.

Slide gốc nêu \textit{Score} của IOI 2001 như một ví dụ: thay vì chỉ hỏi ai
thắng, bài toán cần đánh giá lợi ích bằng điểm số, nên phải dùng quy tắc suy
diễn có mục tiêu hơn.

\subsection{Cài đặt}

Có hai cách cài đặt tự nhiên.

\paragraph{Quy hoạch động.}
Ta liệt kê các trạng thái \((N,K)\) theo thứ tự bảo đảm mọi trạng thái con đã
được tính trước khi tính trạng thái hiện tại. Vì mỗi nước đi làm số sỏi giảm,
đồ thị trạng thái không có chu trình.

\paragraph{Tìm kiếm có nhớ.}
Ta viết hàm đệ quy tính \(f(N,K)\). Khi một trạng thái đã được tính, lưu kết
quả lại để những nhánh khác dùng lại.

Một phiên bản mã giả là:
\[
\begin{array}{l}
\textsc{Solve}(N,K):\\
\quad \textbf{if }N=0\textbf{ then return }\Lose\\
\quad \textbf{if }(N,K)\text{ đã có trong bảng nhớ thì trả lại giá trị đó}\\
\quad \textbf{for }x=1\textbf{ to }\min(N,K)\textbf{ do}\\
\quad\quad \textbf{if }\textsc{Solve}(N-x,2x)=\Lose\textbf{ then}\\
\quad\quad\quad \text{ghi }(N,K)=\Win\text{ và trả về }\Win\\
\quad \text{ghi }(N,K)=\Lose\text{ và trả về }\Lose.
\end{array}
\]

Khi đối cục thực sự diễn ra, chiến lược cũng rất trực tiếp: dựa vào bảng
thắng thua đã biết, luôn chọn một nước đi đưa đối thủ vào trạng thái thua.

\subsection{Độ phức tạp và hạn chế}

Với bài lấy sỏi có \(N\) viên ban đầu, số trạng thái \((n,k)\) cần xét là
\(\Oh(N^2)\). Mỗi trạng thái có thể thử đến \(\Oh(N)\) nước đi, nên cách làm
trực tiếp có độ phức tạp thời gian \(\Oh(N^3)\) và bộ nhớ \(\Oh(N^2)\).

Hạn chế của phương pháp tổng quát đến từ chính nền tảng của nó:
\begin{itemize}
  \item nó dựa trên cấu trúc topo của chuyển trạng thái;
  \item nó phải xét mọi trạng thái liên quan;
  \item nếu số trạng thái quá lớn, thậm chí vô hạn, hoặc quan hệ chuyển trạng
  thái quá phức tạp, thời gian và bộ nhớ đều có thể không chấp nhận được.
\end{itemize}

Vì vậy, phương pháp tổng quát không giải được mọi bài toán đối sách. Để vượt
qua giới hạn đó, ta cần những lời giải riêng cho từng bài: các phương pháp
đặc thù.

\section{Hướng thứ hai: phương pháp đặc thù}

Phương pháp đặc thù không cố dựng toàn bộ đồ thị trạng thái. Thay vào đó, nó
tìm một quy luật ra quyết định để từ trạng thái hiện tại chọn ngay nước đi
tốt. Cách nghĩ này thường bắt đầu từ những cục diện kết thúc hoặc tàn cuộc
nhỏ, rồi phân tích ngược từ dưới lên.

\subsection{Ví dụ đồng xu trên bàn tròn}

Xét trò chơi sau. Trên một mặt bàn tròn, Giáp và Ất luân phiên đặt đồng xu
5 xu, các đồng xu không được chồng lên nhau. Giáp đặt trước; ai đầu tiên
không đặt được nữa thì thua.

Chiến lược thắng của Giáp là đặt đồng xu đầu tiên đúng tại tâm bàn. Sau đó,
mỗi khi Ất đặt một đồng xu ở đâu, Giáp đặt đồng xu đối xứng qua tâm. Vì mặt
bàn tròn đối xứng tâm, nếu Ất đặt hợp lệ thì vị trí đối xứng cũng hợp lệ.
Giáp luôn khôi phục được tính đối xứng và cuối cùng Ất là người hết nước đi.

Trong ví dụ này, Giáp đã tìm được một loại trạng thái thắng mang tính cân
bằng. Mỗi khi đối thủ phá cân bằng, Giáp lập tức khôi phục cân bằng. Câu hỏi
đối với các bài đối sách là: trạng thái cân bằng ở đâu, và quy luật nào giúp
ta khôi phục cân bằng sau khi nó bị phá?

\subsection{Quan sát từ các trường hợp nhỏ}

Trở lại bài lấy sỏi. Nếu tính các giá trị nhỏ từ tàn cuộc, ta nhận được:
\[
\begin{array}{c|cccccccccccc}
N & 2&3&4&5&6&7&8&9&10&11&12&13\\ \hline
f(N,N-1) & \Lose&\Lose&\Win&\Lose&\Win&\Win&
\Lose&\Win&\Win&\Win&\Win&\Lose
\end{array}
\]
Các trạng thái thua đầu tiên của người đi trước là
\[
  2,\ 3,\ 5,\ 8,\ 13,\ldots
\]
Đây chính là dãy Fibonacci với
\[
  F_1=2,\qquad F_2=3,\qquad F_i=F_{i-1}+F_{i-2}\quad(i\ge 3).
\]
Từ đó có phỏng đoán:
\[
  N\in\{F_i\}\Longrightarrow f(N,N-1)=\Lose,
  \qquad
  N\notin\{F_i\}\Longrightarrow f(N,N-1)=\Win.
\]

\subsection{Bốn tính chất dùng trong chứng minh}

Ta ghi lại bốn tính chất cơ bản xuất hiện trong slide.

\paragraph{Tính chất 1.}
Nếu \(K\ge N\), thì trạng thái \((N,K)\) là trạng thái thắng. Người sắp đi
lấy toàn bộ \(N\) viên.

\paragraph{Tính chất 2.}
Nếu \((N,N-1)\) là trạng thái thua, thì với mọi \(K<N\), trạng thái
\((N,K)\) cũng là trạng thái thua. Lý do là tập nước đi từ \((N,K)\) chỉ là
tập con của tập nước đi từ \((N,N-1)\); nếu mọi nước đi trong tập lớn đều dẫn
tới trạng thái thắng cho đối thủ, thì tập nhỏ cũng vậy.

\paragraph{Tính chất 3.}
Nếu \(K<N\) ở trạng thái \((N,K)\), thì trong ván chơi tiếp theo, lượt cuối
cùng lấy hết số sỏi còn lại không thể lấy quá \(2N/3\) viên.

Thật vậy, giả sử lượt cuối cùng lấy \(y\) viên. Ngay trước lượt đó, đối thủ
vừa lấy một số \(z\) sao cho \(y\le 2z\), tức \(z\ge y/2\). Tổng số sỏi từng
có trước hai lượt này ít nhất là \(z+y\ge 3y/2\), mà không vượt quá \(N\).
Do đó \(y\le 2N/3\).

\paragraph{Tính chất 4.}
Với dãy Fibonacci trên,
\[
  \frac{4F_{i-1}}{3}<F_i
\]
trong những bước quy nạp được dùng dưới đây. Với các chỉ số đầu, điều này
kiểm tra trực tiếp; về sau nó theo từ quan hệ truy hồi và tốc độ tăng của
dãy Fibonacci.

\subsection{Kết luận 1: số Fibonacci là trạng thái cân bằng thua}

\paragraph{Kết luận 1.}
Với mọi \(i\), trạng thái \((F_i,A)\) là trạng thái thua nếu \(A<F_i\).

\paragraph{Chứng minh.}
Với \(F_1=2\) và \(F_2=3\), kết luận kiểm tra trực tiếp được.

Giả sử kết luận đúng từ \(F_1\) đến \(F_i\). Ta chứng minh cho
\(F_{i+1}=F_i+F_{i-1}\). Xét một nước đi đầu tiên lấy \(x\) viên.

Nếu \(x\ge F_{i-1}\), người sau nhận trạng thái
\[
  (F_{i+1}-x,2x).
\]
Khi đó
\[
  2x\ge 2F_{i-1}\ge F_{i-1}+F_{i-2}=F_i>F_{i+1}-x,
\]
nên theo tính chất 1, người sau có thể lấy hết phần còn lại và thắng. Vì vậy
nước đi này làm người đi trước thua.

Nếu \(x<F_{i-1}\), xét riêng phần \(F_{i-1}\) trong phân tích
\[
  F_{i+1}=F_i+F_{i-1}.
\]
Theo giả thiết quy nạp, trạng thái \((F_{i-1},x)\) là trạng thái thua. Do đó
người sau có thể chơi trên phần này sao cho sau một số lượt, người đi trước
phải đối mặt với trạng thái có dạng
\[
  (F_i,X).
\]
Theo tính chất 3, trong phần \(F_{i-1}\), lượt cuối cùng lấy không quá
\(2F_{i-1}/3\) viên; do luật nhân đôi, ta có
\[
  X\le 2\cdot \frac{2F_{i-1}}{3}
   =\frac{4F_{i-1}}{3}<F_i.
\]
Vì vậy \((F_i,X)\) là trạng thái thua theo giả thiết quy nạp.

Trong cả hai trường hợp, mọi nước đi đầu tiên từ \((F_{i+1},A)\) đều để lại
cho người sau một cách thắng. Suy ra \((F_{i+1},A)\) là trạng thái thua. Theo
quy nạp, kết luận 1 đúng.

\subsection{Kết luận 2: ngoài Fibonacci là trạng thái thắng}

\paragraph{Kết luận 2.}
Với \(N>2\), nếu \(N\notin\{F_i\}\) thì trạng thái ban đầu \((N,N-1)\) là
trạng thái thắng.

Quy luật quyết định có thể mô tả bằng cách lặp sau. Từ \(N\), tìm số
Fibonacci lớn nhất nhỏ hơn nó rồi trừ đi; với phần dư nhận được, lại tìm số
Fibonacci lớn nhất nhỏ hơn phần dư rồi trừ tiếp. Lặp như vậy cho tới khi phần
dư chính nó là một số Fibonacci. Số còn lại ở bước cuối là số sỏi nên lấy ở
nước đầu.

Nói cách khác, nếu biểu diễn \(N\) thành tổng các số Fibonacci không kề nhau
theo quá trình trên,
\[
  N=F_{i_1}+F_{i_2}+\cdots+F_{i_t}
  \quad(i_1>i_2>\cdots>i_t),
\]
thì người đi trước lấy số hạng nhỏ nhất \(F_{i_t}\). Nước đi này đưa đối thủ
vào một trạng thái cân bằng theo nghĩa của kết luận 1 và các bước quy nạp
trên. Sau đó, mỗi khi đối thủ phá cân bằng, ta dùng cùng quy luật để thu hẹp
phần còn lại và khôi phục cân bằng.

Ví dụ:
\[
\begin{array}{c|c|c}
N & \text{phân tích Fibonacci} & \text{nước đầu thắng}\\ \hline
4 & 3+1 & 1\\
7 & 5+2 & 2\\
12 & 8+3+1 & 1\\
18 & 13+5 & 5
\end{array}
\]
Các dòng này khớp với kết quả tính bằng phương pháp tổng quát.

\subsection{Ý nghĩa thuật toán}

Trong bài lấy sỏi, trạng thái cân bằng chính là các số Fibonacci. Quy luật
quyết định là liên tục thu hẹp phạm vi bằng cách tìm số Fibonacci lớn nhất
phù hợp, cho tới khi xác định được số sỏi cần lấy.

So với phương pháp tổng quát, độ phức tạp giảm mạnh:
\[
\begin{array}{c|c|c}
\text{Phương pháp} & \text{Bộ nhớ} & \text{Thời gian}\\ \hline
\text{Tổng quát} & \Oh(N^2) & \Oh(N^3)\\
\text{Đặc thù} & \Oh(1) & \Oh(\log N)
\end{array}
\]
Ở đây \(\Oh(\log N)\) đến từ việc sinh hoặc tra các số Fibonacci tới \(N\).
Nếu cần cài đặt cho số lớn, có thể lưu vài số Fibonacci liên tiếp và dịch dần
theo quan hệ truy hồi, nên bộ nhớ vẫn là hằng số.

\section{Khung của phương pháp đặc thù}

Phương pháp đặc thù vẫn phải mô tả trạng thái: ta liệt kê những yếu tố ảnh
hưởng đến kết quả thắng thua và tổng hợp chúng thành trạng thái. Điểm khác
là ta không cần dựng toàn bộ cấu trúc chuyển trạng thái.

Thay vào đó, ta phân tích ngược từ cục diện kết thúc hoặc những tàn cuộc nhỏ.
Trong các trường hợp đặc biệt như quy mô nhỏ, đối xứng, giá trị cực đại, giá
trị cực tiểu, ta tìm những lớp trạng thái thắng hoặc thua của người đi trước,
rồi quan sát điểm chung của chúng. Những điểm chung thường rơi vào ba hướng:
\begin{itemize}
  \item tính đối xứng;
  \item tính đơn giản của mô tả;
  \item tính dị biệt so với các trạng thái lân cận.
\end{itemize}

Sau khi tìm được tính chất chung, ta mở rộng nó sang trường hợp tổng quát.
Kết quả mong muốn là một lớp trạng thái cân bằng, thường là trạng thái thắng
hoặc thua chắc, cùng một quy luật ra quyết định giúp duy trì cân bằng đó.

Trong slide gốc, các ví dụ được nêu là bài \textit{Polygon} của POI 1999 và
bài lấy số của IOI 1996: chúng đều có thể được tiếp cận bằng hướng đặc thù.

\section{So sánh hai hướng}

Nếu mở rộng bài lấy sỏi bằng cách cho mỗi lượt lấy nhiều nhất \(3\), \(4\),
\(5\), hoặc nói chung \(k\) lần số sỏi đối thủ vừa lấy, ta lại đứng trước lựa
chọn giữa hai hướng:
\[
  \text{phương pháp tổng quát}
  \qquad\text{và}\qquad
  \text{phương pháp đặc thù}.
\]

\subsection{Chiều phân tích}

Phương pháp tổng quát đi từ trên xuống: bắt đầu từ trạng thái ban đầu và xét
thắng thua của mọi trạng thái liên quan.

Phương pháp đặc thù đi từ dưới lên: bắt đầu từ kết thúc hoặc tàn cuộc nhỏ,
rồi nghiên cứu một lớp trạng thái cân bằng.

\subsection{Quy tắc thắng thua}

Phương pháp tổng quát có quy tắc phổ dụng:
\[
\begin{array}{l}
\text{có một nước đi tới trạng thái thua}\Rightarrow \text{thắng},\\
\text{mọi nước đi đều tới trạng thái thắng}\Rightarrow \text{thua}.
\end{array}
\]

Phương pháp đặc thù không có quy tắc phổ dụng. Mỗi bài cần tìm cấu trúc riêng
của nó.

\subsection{Cài đặt}

Với phương pháp tổng quát, khâu quan trọng là tiền xử lý bằng quy hoạch động
hoặc tìm kiếm có nhớ.

Với phương pháp đặc thù, khâu quan trọng là suy nghĩ trước để phát hiện quy
luật quyết định, rồi chuyển quy luật đó thành chương trình.

\subsection{Ưu và nhược điểm}

\begin{center}
\begin{tabular}{p{0.24\linewidth}|p{0.34\linewidth}|p{0.34\linewidth}}
\textbf{Phương pháp} & \textbf{Ưu điểm} & \textbf{Nhược điểm}\\ \hline
Tổng quát &
Có quy tắc chung, áp dụng rộng. &
Bị giới hạn bởi đồ thị chuyển trạng thái; phải xét nhiều trạng thái, có thể
tốn thời gian và bộ nhớ.\\[0.4em]
Đặc thù &
Không phụ thuộc vào việc dựng toàn bộ đồ thị; thường nhanh, tiết kiệm bộ nhớ
và dễ cài đặt sau khi đã có quy luật. &
Không có mẫu chung; phần khó nằm ở phân tích và phát hiện quy luật.
\end{tabular}
\end{center}

\subsection{Tư tưởng cốt lõi}

Trong bài toán đối sách loại này, một trạng thái hoặc là thắng chắc cho người
sắp đi, hoặc là thua chắc cho người sắp đi. Vì vậy khi chơi, mục tiêu của ta
là chiếm trạng thái thắng và để lại trạng thái thua cho đối phương.

Đó là tư tưởng trung tâm của việc giải bài toán đối sách.

Nhìn rộng hơn, phương pháp tổng quát xét mọi trạng thái từ một góc nhìn thống
nhất để quyết định chiến lược; nó tương tự tinh thần của quy hoạch động.
Phương pháp đặc thù xét một lớp trạng thái quan trọng từ góc nhìn riêng; nó
gần với tinh thần tham lam, vì sau khi tìm được quy luật, mỗi bước đều chọn
ngay quyết định giữ được cấu trúc mong muốn.

\section{Kết luận}

Lý thuyết đối sách là một nhánh quan trọng của vận trù học. Thông qua bài lấy
sỏi, bài viết đã trình bày hai hướng giải một lớp bài toán đối sách: phương
pháp tổng quát và phương pháp đặc thù. Hai hướng này vừa là phương pháp, vừa
là cách tư duy.

Kinh nghiệm trên giấy luôn cần được kiểm chứng bằng thực hành. Khi luyện tập
nhiều, ta sẽ biết phối hợp tư duy xuôi và tư duy ngược, quan sát bài toán từ
nhiều góc độ, đi từ tổng quát đến đặc biệt rồi từ đặc biệt quay lại tổng
quát, và chọn cách cài đặt phù hợp cho từng bài.

Một cách nói cô đọng là: lập mưu trong màn trướng, quyết thắng ngoài nghìn
dặm.

\end{document}
